\paragraph{De la codificación ternaria a la bandera marcada.}
\label{inc-ampl-bandera}
Las regiones aportan también los códigos $w_\pi,w_e,w_\varphi$ y
$w_{\rm APP}$ en el marco heredado. La aplicación
\[
 \gamma_{A_W}:w\longmapsto(w,wA_W)
\]
conserva la palabra de seis componentes como primera coordenada y añade
su transformación por $A_W$. La relación $A_W^2=-I_6$ fija la
compatibilidad cuadrática de esta realización. El paso siguiente,
$c\mapsto\operatorname{supp}(c)$, conserva las posiciones no nulas.
En las palabras de peso seis, el soporte identifica exactamente la pareja
$\{c,-c\}$, como prueba la distancia mínima del código. De esta forma,
el cociente por signo produce las hexadas, mientras el código completo
sigue disponible para las operaciones que requieren sus símbolos.

En esta coordenatización, las dos construcciones de la cuaterna coinciden:
\begin{equation}
 \{i:K_i\geq729\}
 =H_e\cap P_\pi=B_K=\{4,6,9,10\}.
 \label{inc-ampl-cuaterna}
\end{equation}
El lado izquierdo procede de una comparación entera de los componentes
de $K$; el derecho, de la intersección de los soportes codificados de
propagación y clausura. El umbral pertenece a la coordenatización de bloques
declarada y permanece explicitado en la definición del lector. El
registro íntegro conserva los valores sobre los que actúa, incluso cuando
distintos umbrales producen el mismo subconjunto. La igualdad expresa,
por tanto, una compatibilidad entre aplicaciones especificadas.

El soporte $H^-_\alpha$ completa esa cuaterna hasta una hexada. La
condición de pertenencia al soporte $P_\pi$, junto con los tamaños de
intersección $(4,3,2)$ respecto de $H_e,H_\varphi,H_{\rm APP}$,
selecciona una única hexada. Las pruebas siguientes reúnen las dos
determinaciones: la obtenida mediante el signo de $Z_0$ y la obtenida
mediante incidencia. El resultado se organiza como
\begin{equation}
 \widetilde{\mathcal I}_*
 =\bigl(B_K\subset H^-_\alpha;
         P_\pi,H_e,H_\varphi,H_{\rm APP}\bigr).
 \label{inc-ampl-marco}
\end{equation}
La inclusión es la bandera; los cuatro soportes adicionales constituyen
su marco marcado. Las permutaciones simultáneas transportan el marco y
conservan sus inclusiones e intersecciones. La información relevante reside
en esta configuración transportada y en su clase de isomorfía.

Las cinco regiones de $\pi$ recuperan aquí su diferenciación interna.
Comparten $w_\pi$ como código ternario visible, pero conservan emisiones,
firmas y multiplicidades distintas. En la realización binaria sobre una
dodecada marcada, la región transversal corresponde a la octada cuya
intersección con la dodecada es la cuaterna; las otras cuatro regiones
corresponden a las elevaciones de las cuatro hexadas de su estrella. La
región negativa queda distinguida por $H^-_\alpha$ y el orden de coordenatización
individualiza las tres regiones positivas. Esta correspondencia conserva
los papeles regionales y la distribución $432+4\cdot144=1008$, además
de realizar la partición incidencial $1+4$.
